\section{Retained reconstruction statements and resource bounds}
\label{sec:localized-statements}
\label{sec:retained-summaries}
We retain the localized experiment, its binary converse and the earlier
stationary bounds. Their proofs follow in the remaining appendices and
in Sections~\ref{sec:stationary}--\ref{sec:stationary-information}.
The stronger stationary bounds proved in Section~\ref{sec:rare-stationary}
use the additional prior-fixed compass design stated there.

\begin{theorem}[Adaptive scalar reconstruction]\label{thm:adaptive}
Fix the numerical bounds of $\mathcal B_\eta$ and the kernel above.
There are $C,c,\nu_0>0$ such that, for $0<\nu<\nu_0$ and
$0<\delta<1/4$, a finite adaptive experiment with the reciprocal
compass commands reconstructs a smooth strictly convex periodic
presentation with geometric loss at most $C\nu$, with probability at
least $1-\delta$. It recovers the true primitive orbit count and exact
rational relations of a true period-generating list in a true independent-pair basis. The Euclidean coordinates of these vectors are estimates.

After a prior-dependent coarse acquisition, the finest localization
scale and sufficient coupled calibration are
\begin{equation}\label{eq:precision}
 \sigma=c\nu^{s/(s-2)},\qquad \ell+\tau+t\alpha\le c\sigma.
\end{equation}
The number of nominal command centers, counted with repetitions if
necessary, and the number of attempted preparations satisfy
\begin{align}
 J_\nu&\le C\nu^{-1/(s-2)}\log(C/\nu),\label{eq:centers}\\
 N_\nu&\le C\nu^{-1/(s-2)}\log(C/\nu)
                               \log(C/(\nu\delta)).\label{eq:attempts}
\end{align}
Every solid start and free miss is charged in $N_\nu$. All commands lie
in one fixed prior-controlled aperture. The reconstruction uses finite
local computations, radial bisections and finite polynomial data, not
a search over an infinite-dimensional class of physical candidates.
Its constants include the exit-time and patch-margin bounds.
\end{theorem}

Here and below a finite smooth partition of unity, with its derivatives,
is part of the fixed numerical representation. Certified evaluations
to the reserved tolerances suffice. The theorem counts preparations and
centers, not apparatus travel or bit operations for elementary functions.
Known $s$ is used; adaptivity refers to the locations of commands, not
adaptation to unknown smoothness.

\begin{theorem}[A necessary number of binary observations]
\label{thm:lower}
For each fixed $0<\beta\le1$ there is a fixed choice of the bounds
above and a nondegenerate subclass $\mathcal P\subset\mathcal B_\eta$
with the following property. Any adaptive randomized procedure whose
only table-dependent outputs are at most $N$ bits, and which recovers
the component boundaries with $C^2$ error at most $\nu$ with probability
at least $3/4$ uniformly on $\mathcal P$, must satisfy
\begin{equation}\label{eq:lower}
                         N\ge c\nu^{-1/(s-2)}.
\end{equation}
This holds even when the lattice, its primitive orbit count, component
correspondences and laboratory pose are given, and implementation is
exact. Input choices may have arbitrary precision and depend on all
previous bits. Independent controller randomness is allowed.
\end{theorem}

The construction of $\mathcal P$ is given in Section~\ref{sec:lower}.
The assertion is not a lower bound for every subclass, in particular
not for a singleton or a finite parametric family. On any bounded class containing this construction, $\eqref{eq:attempts}$ and
$\eqref{eq:lower}$ identify the power of the attempted-bit complexity.
A logarithmic gap remains. Neither result prices input precision or
converts an active experiment into a passive billiard invariant.
